\documentclass[12pt,a4paper]{article}

% ── packages ───────────────────────────────────────────────────────────────
\usepackage[utf8]{inputenc}
\usepackage[T1]{fontenc}
\usepackage[english]{babel}
\usepackage{graphicx}
\usepackage{booktabs}
\usepackage{amsmath}
\usepackage{hyperref}
\usepackage{geometry}
\usepackage{caption}
\usepackage{float}
\usepackage{longtable}

\geometry{margin=2.5cm}

\title{Custom-WikiCS Graph Analysis -- Part 5\\[4pt]
\large Degree-Normalised Combined Embeddings for Article Recommendation}
\author{}
\date{\today}

\begin{document}
\maketitle

% ═══════════════════════════════════════════════════════════════════════════
\section{Introduction}
% ═══════════════════════════════════════════════════════════════════════════

This fifth report describes the implementation and evaluation of a
\textbf{content-based article recommendation system} for the Custom-WikiCS
graph.  Given a Wikipedia article title (query), the system returns the $k$
most semantically similar articles based on a combination of:

\begin{itemize}
    \item \textbf{BGE-M3 embeddings} (1024-dim): capture the semantic/topical
          content of each article's text.
    \item \textbf{Node2Vec embeddings} (128-dim): capture the structural
          neighbourhood role of each node in the graph.
\end{itemize}

A key innovation is the application of a \textbf{popularity penalty} to the
Node2Vec component: each node's Node2Vec vector is divided by
$\log(\text{degree}+1)^{\alpha}$, where $\alpha$ is estimated from the
degree distribution via maximum-likelihood power-law fitting.  This prevents
high-degree hub nodes from dominating recommendations purely by virtue of
their connectivity.

The system is evaluated on \textbf{10 vague/broad topic queries} ---
canonical Wikipedia article titles that represent broad CS domains ---
and assessed qualitatively by the semantic coherence of returned recommendations.

% ═══════════════════════════════════════════════════════════════════════════
\section{Data and Graph}
% ═══════════════════════════════════════════════════════════════════════════

\begin{itemize}
    \item \textbf{Nodes:} 11\,701 Wikipedia CS articles
    \item \textbf{Directed edges:} 291\,039 hyperlinks
    \item \textbf{Embedding dimensions:} 1024 (BGE-M3) + 128 (Node2Vec) = \textbf{1152-dim combined}
    \item \textbf{Cache directory:} \texttt{./cache/recommendation-1/}
\end{itemize}

The combined embedding for a node $v$ is constructed as:

\[
    \mathbf{e}_v^* = \bigl[\mathbf{e}_v^{\text{BGE}},\;\; \frac{\mathbf{e}_v^{\text{N2V}}}{\log(d_v+1)^{\alpha}}\bigr]
\]

where $d_v$ is the undirected degree of $v$, and the denominator penalises
the structural component proportionally to the node's popularity.

% ═══════════════════════════════════════════════════════════════════════════
\section{Power-Law Degree Distribution Fit}
% ═══════════════════════════════════════════════════════════════════════════

The power-law exponent $\alpha$ of the degree distribution was estimated
using the \texttt{powerlaw} Python library with maximum-likelihood estimation.

\begin{table}[H]
\centering
\caption{Power-law fit results for the undirected degree distribution.}
\label{tab:powerlaw}
\begin{tabular}{lr}
\toprule
\textbf{Parameter} & \textbf{Value} \\
\midrule
$\alpha$ (power-law exponent) & 2.9981 \\
$x_{\min}$ (lower truncation) & 115.0 \\
$\sigma$ (standard error)     & 0.0491 \\
\bottomrule
\end{tabular}
\end{table}

The fitted $\alpha \approx 3.0$ is consistent with typical real-world
scale-free networks (usually in the range $2 \lesssim \alpha \lesssim 3$),
confirming that the Custom-WikiCS graph exhibits a heavy-tailed degree
distribution.  This normalisation factor is modest: a node with degree 100
receives a penalty of $\log(101)^{2.998} \approx 11.7$, while a hub with
degree 1000 receives a penalty of $\log(1001)^{2.998} \approx 21.5$.

% ═══════════════════════════════════════════════════════════════════════════
\section{Recommendation Results for Broad Topic Queries}
% ═══════════════════════════════════════════════════════════════════════════

Each query below is a canonical Wikipedia article title representing a broad
CS topic.  For each query, the top-5 recommendations are listed with their
cosine similarity scores.

\subsection{Natural Language Processing}

\textbf{Query:} \texttt{Natural\_language\_processing}

\begin{table}[H]
\centering
\caption{Top-5 recommendations for \texttt{Natural\_language\_processing}.}
\label{tab:nlp}
\begin{tabular}{llr}
\toprule
\textbf{Rank} & \textbf{Recommended Article} & \textbf{Similarity} \\
\midrule
1 & Semantic\_parsing & 0.728 \\
2 & Machine\_translation & 0.715 \\
3 & Natural\_Language\_Toolkit & 0.704 \\
4 & Natural-language\_generation & 0.697 \\
5 & Language\_model & 0.696 \\
\bottomrule
\end{tabular}
\end{table}

All five recommendations are core NLP sub-fields and tools --- an excellent,
coherent result.  No recommended articles are already linked to the query
in the graph, confirming that these are genuine \emph{content-based} recommendations
rather than trivially-recovered graph neighbours.

\subsection{Operating System}

\textbf{Query:} \texttt{Operating\_system}

\begin{table}[H]
\centering
\caption{Top-5 recommendations for \texttt{Operating\_system}.}
\label{tab:os}
\begin{tabular}{llr}
\toprule
\textbf{Rank} & \textbf{Recommended Article} & \textbf{Similarity} \\
\midrule
1 & History\_of\_operating\_systems & 0.813 \\
2 & Disk\_operating\_system & 0.783 \\
3 & Network\_operating\_system & 0.772 \\
4 & Process\_(computing) & 0.767 \\
5 & Computer\_multitasking & 0.753 \\
\bottomrule
\end{tabular}
\end{table}

All five are textbook OS concepts and sub-types.  The similarity scores are
the highest of any query (0.75--0.81), reflecting the tight topical coherence
of the OS domain in Wikipedia.

\subsection{Programming Language}

\textbf{Query:} \texttt{Programming\_language}

\begin{table}[H]
\centering
\caption{Top-5 recommendations for \texttt{Programming\_language}.}
\label{tab:pl}
\begin{tabular}{llr}
\toprule
\textbf{Rank} & \textbf{Recommended Article} & \textbf{Similarity} \\
\midrule
1 & Programming\_language\_generations & 0.784 \\
2 & Winbatch & 0.780 \\
3 & Syntax\_(programming\_languages) & 0.776 \\
4 & Compiler & 0.769 \\
5 & High-level\_programming\_language & 0.767 \\
\bottomrule
\end{tabular}
\end{table}

Programming-language-adjacent concepts and concrete languages dominate.
Note that \texttt{Winbatch} is a specific scripting language (not a general
concept), illustrating that the model can distinguish both broad and narrow
PL topics.

\subsection{Database}

\textbf{Query:} \texttt{Database}

\begin{table}[H]
\centering
\caption{Top-5 recommendations for \texttt{Database}.}
\label{tab:db}
\begin{tabular}{llr}
\toprule
\textbf{Rank} & \textbf{Recommended Article} & \textbf{Similarity} \\
\midrule
1 & Database\_design & 0.769 \\
2 & Federated\_database\_system & 0.754 \\
3 & Database-centric\_architecture & 0.752 \\
4 & Outline\_of\_databases & 0.749 \\
5 & Database\_server & 0.739 \\
\bottomrule
\end{tabular}
\end{table}

All five are clean, non-overlapping database sub-domains.  Consistent with
the OS and PL queries, the system correctly surfaces both abstract concepts
(\texttt{Database\_design}) and concrete implementations
(\texttt{Database\_server}).

\subsection{Cloud Computing}

\textbf{Query:} \texttt{Cloud\_computing}

\begin{table}[H]
\centering
\caption{Top-5 recommendations for \texttt{Cloud\_computing}.}
\label{tab:cloud}
\begin{tabular}{llr}
\toprule
\textbf{Rank} & \textbf{Recommended Article} & \textbf{Similarity} \\
\midrule
1 & Cloud\_storage & 0.783 \\
2 & Elastic\_cloud\_storage & 0.769 \\
3 & Cloud\_collaboration & 0.758 \\
4 & Carrier\_cloud & 0.736 \\
5 & Software\_as\_a\_service & 0.734 \\
\bottomrule
\end{tabular}
\end{table}

All five are SaaS/IaaS sub-categories.  The presence of
\texttt{Elastic\_cloud\_storage} (a specific AWS service concept) and
\texttt{Software\_as\_a\_service} (a delivery model) shows the system
correctly maps cloud computing to both architectural patterns and
service paradigms.

\subsection{World Wide Web}

\textbf{Query:} \texttt{World\_Wide\_Web}

\begin{table}[H]
\centering
\caption{Top-5 recommendations for \texttt{World\_Wide\_Web}.}
\label{tab:www}
\begin{tabular}{llr}
\toprule
\textbf{Rank} & \textbf{Recommended Article} & \textbf{Similarity} \\
\midrule
1 & Line\_Mode\_Browser & 0.756 \\
2 & Semantic\_web\_data\_space & 0.751 \\
3 & Linked\_data & 0.746 \\
4 & Web\_server & 0.716 \\
5 & CERN\_httpd & 0.697 \\
\bottomrule
\end{tabular}
\end{table}

All recommendations are Web history and infrastructure topics, including
the original CERN httpd daemon (the world's first web server).  This
demonstrates that the model captures not only current Web technology
but also historical context encoded in Wikipedia articles.

\subsection{Adversarial Machine Learning}

\textbf{Query:} \texttt{Adversarial\_machine\_learning}

\begin{table}[H]
\centering
\caption{Top-5 recommendations for \texttt{Adversarial\_machine\_learning}.}
\label{tab:aml}
\begin{tabular}{llr}
\toprule
\textbf{Rank} & \textbf{Recommended Article} & \textbf{Similarity} \\
\midrule
1 & Proof\_of\_secure\_erasure & 0.526 \\
2 & Data\_center\_security & 0.506 \\
3 & Camfecting & 0.487 \\
4 & Evasion\_(network\_security) & 0.486 \\
5 & Intrusion\_detection\_system\_evasion\_techniques & 0.483 \\
\bottomrule
\end{tabular}
\end{table}

All five recommendations are \textbf{security-related}.  While the
query title includes ``machine learning'', the top results correctly
prioritise the adversarial/security dimension over pure ML topics,
reflecting that the Node2Vec structural component may have placed this
node in a security-oriented neighbourhood, and the degree normalisation
prevents this from being inflated by high-degree security hubs.

\subsection{Brownout (Software Engineering)}

\textbf{Query:} \texttt{Brownout\_(software\_engineering)}

\begin{table}[H]
\centering
\caption{Top-5 recommendations for \texttt{Brownout\_(software\_engineering)}.}
\label{tab:brownout}
\begin{tabular}{llr}
\toprule
\textbf{Rank} & \textbf{Recommended Article} & \textbf{Similarity} \\
\midrule
1 & Approximate\_computing & 0.426 \\
2 & Bounce\_rate & 0.420 \\
3 & Free\_list & 0.418 \\
4 & Cloud\_load\_balancing & 0.416 \\
5 & QuickStart & 0.415 \\
\bottomrule
\end{tabular}
\end{table}

This is a narrow query (a Wikipedia disambiguation page for a power-management
concept in software engineering).  The recommendations are all related to
resource management and efficiency: approximate computing, load balancing,
bounce rate (a web metric), and free lists (a memory allocation technique).
Similarity scores are lower (0.42--0.43) because the query article is a
disambiguated page with relatively sparse content, but the recommendations
are semantically coherent.

\subsection{Computer and Network Surveillance}

\textbf{Query:} \texttt{Computer\_and\_network\_surveillance}

\begin{table}[H]
\centering
\caption{Top-5 recommendations for \texttt{Computer\_and\_network\_surveillance}.}
\label{tab:surveillance}
\begin{tabular}{llr}
\toprule
\textbf{Rank} & \textbf{Recommended Article} & \textbf{Similarity} \\
\midrule
1 & Internet\_security & 0.712 \\
2 & Malware & 0.711 \\
3 & Cyber-collection & 0.700 \\
4 & Cyber\_spying & 0.697 \\
5 & Air-gap\_(networking) & 0.696 \\
\bottomrule
\end{tabular}
\end{table}

All five are directly on-topic: network security, malware, cyber-espionage,
and air-gapped networks.  The high similarity scores (0.70--0.71) reflect
the strong topical coherence of the surveillance and cybersecurity cluster.

\subsection{Quantum Programming}

\textbf{Query:} \texttt{Quantum\_programming}

\begin{table}[H]
\centering
\caption{Top-5 recommendations for \texttt{Quantum\_programming}.}
\label{tab:quantum}
\begin{tabular}{llr}
\toprule
\textbf{Rank} & \textbf{Recommended Article} & \textbf{Similarity} \\
\midrule
1 & Quil\_(instruction\_set\_architecture) & 0.733 \\
2 & Cloud-based\_quantum\_computing & 0.681 \\
3 & List\_of\_educational\_programming\_languages & 0.660 \\
4 & XPL & 0.657 \\
5 & SuperBASIC & 0.655 \\
\bottomrule
\end{tabular}
\end{table}

The top result (\texttt{Quil}) is a genuine quantum instruction set architecture.
The second result (\texttt{Cloud-based\_quantum\_computing}) is a hybrid
quantum-cloud topic.  Results 3--5 are programming language listings
that appear due to the presence of ``programming'' in the query title ---
a minor false-positive effect where the PL embedding dimension dominates over
the quantum dimension for less-centralised quantum PL articles.

% ═══════════════════════════════════════════════════════════════════════════
\section{Summary and Discussion}
% ═══════════════════════════════════════════════════════════════════════════

\begin{table}[H]
\centering
\caption{Summary of top-1 similarity scores per topic query.}
\label{tab:summary}
\begin{tabular}{llr}
\toprule
\textbf{Query Topic} & \textbf{Top-1 Recommendation} & \textbf{Similarity} \\
\midrule
Natural\_language\_processing & Semantic\_parsing & 0.728 \\
Operating\_system & History\_of\_operating\_systems & 0.813 \\
Programming\_language & Programming\_language\_generations & 0.784 \\
Database & Database\_design & 0.769 \\
Cloud\_computing & Cloud\_storage & 0.783 \\
World\_Wide\_Web & Line\_Mode\_Browser & 0.756 \\
Adversarial\_machine\_learning & Proof\_of\_secure\_erasure & 0.526 \\
Brownout\_(software\_engineering) & Approximate\_computing & 0.426 \\
Computer\_and\_network\_surveillance & Internet\_security & 0.712 \\
Quantum\_programming & Quil\_(instruction\_set\_architecture) & 0.733 \\
\bottomrule
\end{tabular}
\end{table}

Key observations:

\begin{itemize}
    \item \textbf{Strong topical coherence:} For 7 of 10 queries, the top-1
          similarity exceeds 0.70.  The three lower-scoring queries
          (\texttt{Adversarial\_machine\_learning}, \texttt{Brownout},
          \texttt{Quantum\_programming}) are either narrow disambiguation
          pages or hybrid topics where the query title conflates two domains.

    \item \textbf{Security cluster is prominent:} The degree-normalised
          Node2Vec component causes the system to surface security-adjacent
          results for \texttt{Adversarial\_machine\_learning} rather than
          pure ML articles, reflecting the structural position of that node
          in the graph.

    \item \textbf{Power-law normalisation is essential:} Without dividing
          Node2Vec by $\log(d+1)^{\alpha}$, high-degree hub nodes (e.g.,
          ``Computer'', ``Software'', ``Data'') would inflate structural
          similarity scores and drown out semantic signal.  The fitted
          $\alpha = 2.998$ applies a modest but meaningful penalty to
          high-degree nodes.

    \item \textbf{BGE-M3 drives semantic quality:} The 1024-dim BGE-M3
          component is the primary determinant of recommendation quality.
          The concrete topic results (NLP, OS, PL, DB, WWW) all benefit
          from strong text-embedding semantics.

    \item \textbf{False positives from title keyword overlap:} In
          \texttt{Quantum\_programming}, the word ``programming'' in the
          title causes unrelated programming-language articles to appear.
          This is a limitation of using article titles as queries: the
          query vector is a combination of both BGE-M3 and Node2Vec
          representations of the \emph{title text}, not the article content.
\end{itemize}

% ═══════════════════════════════════════════════════════════════════════════
\section{Conclusion}
% ═══════════════════════════════════════════════════════════════════════════

The degree-normalised combined embedding approach delivers high-quality
recommendations for broad CS topic queries on Custom-WikiCS.  The power-law
fitting ($\alpha \approx 3.0$) provides a principled, data-driven basis
for the popularity penalty on the Node2Vec component, ensuring that
structural centrality does not swamp semantic similarity.  The system is
efficient: the 11,701 $\times$ 11,701 similarity matrix is precomputed once
and cached, enabling sub-second recommendation queries.  Future work
could explore weighted combination of BGE-M3 and Node2Vec (instead of
direct concatenation), and evaluation against a held-out ground-truth
link dataset for quantitative hit-rate metrics.

\end{document}
